\chapter{Wstęp}

Zarządzanie projektem informatycznym obejmuje organizowanie pracy zespołu, planowanie harmonogramu, kontrolowanie kosztów oraz reagowanie na zagrożenia występujące podczas realizacji zadań. Działania te wymagają uwzględnienia zależności pomiędzy zakresem projektu, dostępnymi zasobami i oczekiwanym terminem zakończenia prac. Zmiana jednego z tych elementów może wpływać na pozostałe, co sprawia, że planowanie wymaga regularnej aktualizacji i oceny przyjętych założeń.

Przedmiotem niniejszej pracy jest system PlanWise, którego zadaniem jest połączenie organizacji pracy projektowej z analizą danych wspierającą podejmowanie decyzji. Rozwiązanie obejmuje zarówno funkcje zarządzania projektami, zadaniami i sprintami, jak i mechanizmy oceny ryzyka, priorytetyzacji backlogu, wspomagania przydziału zadań oraz estymacji kosztów. Praca koncentruje się na zaprojektowaniu i implementacji tych funkcji oraz określeniu sposobu oceny ich przydatności.

\section{Uzasadnienie wyboru tematu}

Podstawowym problemem rozpatrywanym w pracy jest wykorzystanie danych projektowych do wspierania decyzji dotyczących dalszego przebiegu przedsięwzięcia. Rejestrowanie zadań, terminów i osób odpowiedzialnych pozwala uporządkować informacje, jednak samo ich zgromadzenie nie rozstrzyga, które zadania należy wykonać w pierwszej kolejności, gdzie może wystąpić opóźnienie ani jak rozdzielić pracę pomiędzy członków zespołu. Odpowiedź na te pytania wymaga uwzględnienia wielu czynników jednocześnie.

Przykładem jest ustalanie kolejności realizacji backlogu. Zadanie o wysokiej wartości biznesowej może wymagać wcześniejszego wykonania prac technicznych, które samodzielnie nie dostarczają wartości użytkownikowi, a zadanie o niewielkiej pracochłonności może blokować kilka kolejnych elementów projektu. Podobne zależności występują przy przydzielaniu pracy: kierowanie się wyłącznie kompetencjami prowadzi do przeciążenia najlepiej przygotowanych osób, a wyłącznie dostępnością — do przypisania zadania komuś bez odpowiedniego przygotowania. Drugim zagadnieniem jest niepewność dotycząca czasu i kosztów: opisy wymagań bywają na etapie planowania niepełne, a oszacowania zmieniają się wraz z poznawaniem szczegółów rozwiązania, co uzasadnia przedstawianie wariantów realizacji wraz z założeniami, od których zależą.

Zagadnienia te uzasadniają opracowanie systemu łączącego dane operacyjne z metodami analitycznymi. W PlanWise wykorzystano metody heurystyczne, uczenie maszynowe, algorytmy harmonogramowania oraz duże modele językowe, przy czym każde z tych podejść odpowiada innemu rodzajowi problemu: reguły punktowe umożliwiają jawną ocenę czynników ryzyka, model uczony wykorzystuje zależności obecne w danych historycznych, rozwiązanie zadania z ograniczeniami wyznacza przydział zgodny z kolejnością prac, a model językowy pozwala sięgnąć po tekstowe opisy zadań.

Wybór tematu wynika również z możliwości połączenia zagadnień inżynierii oprogramowania z oceną metod wspomagania decyzji. Opracowanie systemu wymaga spójnego modelu danych, określenia odpowiedzialności modułów i przepływu informacji między częścią organizacyjną a analityczną, natomiast część badawcza pozwala sprawdzić, w jakich warunkach generowane wyniki są użyteczne i jakie ograniczenia należy uwzględnić przy ich interpretacji.

\section{Cel pracy}

Głównym celem pracy jest zaprojektowanie, implementacja i ocena systemu informatycznego PlanWise, wspierającego kierownika projektu w organizacji pracy oraz podejmowaniu decyzji dotyczących ryzyka, priorytetów, harmonogramu i kosztów.

Cel inżynierski obejmuje opracowanie aplikacji internetowej z warstwą serwerową w technologii .NET oraz interfejsem użytkownika wykonanym w Angularze. System ma umożliwiać zarządzanie projektami i zespołami, definiowanie zadań i ich zależności, prowadzenie backlogu, obsługę sprintów oraz prezentowanie postępu prac na tablicy Kanban i wykresie Gantta. Funkcje te mają dostarczać uporządkowanych danych dla mechanizmów analitycznych.

W obszarze wspomagania decyzji celem jest opracowanie oceny ryzyka opóźnień, rekomendacji kolejności realizacji zadań oraz propozycji przydziału pracy uwzględniających charakterystykę zespołu. Uzupełnieniem jest estymacja kosztów z wykorzystaniem dużego modelu językowego, obejmująca wariant optymistyczny, realistyczny i pesymistyczny oraz uzasadnienie przyjętych założeń i rekomendacje redukcji kosztów.

Cel badawczy polega na ocenie jakości wyników dostarczanych przez wybrane mechanizmy oraz porównaniu ich z określonymi metodami bazowymi. W odniesieniu do predykcji opóźnień obejmuje on dopasowanie modelu uczenia maszynowego, zmierzenie jego jakości względem metody odniesienia w warunkach odpowiadających rzeczywistemu użyciu oraz określenie, jakie dane są potrzebne, aby porównać go z istniejącą oceną heurystyczną. Dla pozostałych mechanizmów ocena obejmuje zgodność rekomendacji z kryteriami, jakość przydziału pracy mierzoną względem zadeklarowanej funkcji celu oraz poprawność i powtarzalność estymacji kosztów.

Istotnym założeniem jest przedstawianie użytkownikowi przesłanek stojących za rekomendacją. Kierownik projektu powinien mieć możliwość przeanalizowania wyniku, poznania jego ograniczeń oraz podjęcia decyzji o zastosowaniu proponowanych zmian.

\section{Zakres pracy}

Zakres pracy obejmuje analizę potrzeb użytkowników, przegląd wybranych rozwiązań wspierających zarządzanie projektami, określenie wymagań, zaprojektowanie architektury, implementację systemu oraz przygotowanie i przeprowadzenie jego oceny. Głównym obszarem zastosowania są projekty informatyczne, w których praca jest organizowana za pomocą zadań, backlogu i sprintów.

Część organizacyjna systemu obejmuje obsługę tożsamości i dostępu, zarządzanie projektami oraz uczestnictwem w zespołach. Dane członków zespołu uwzględniają role, kompetencje, deklarowaną pojemność i stawki godzinowe. Zarządzanie pracą obejmuje zadania, podzadania, zależności, komentarze, statusy i przypisanie do sprintów. System udostępnia również historię aktywności, powiadomienia oraz aktualizację wybranych danych w czasie rzeczywistym.

Architektura backendu opiera się na modularnym monolicie. Wyróżniono moduły tożsamości i dostępu, zarządzania projektami i zespołami, realizacji zadań, harmonogramowania, oceny ryzyka, priorytetyzacji backlogu, estymacji kosztów oraz powiadomień. Zakres opracowania obejmuje podział odpowiedzialności pomiędzy modułami, organizację warstw, komunikację poprzez wspólne kontrakty, wykorzystanie zdarzeń domenowych i przetwarzanie operacji analitycznych w tle. Dane przechowywane są w relacyjnej bazie PostgreSQL.

Część analityczna obejmuje ocenę ryzyka na podstawie parametrów zadań oraz prognozowanie realizacji sprintu. Rozwiązanie heurystyczne stanowi metodę bazową i punkt odniesienia dla modelu uczenia maszynowego, przy czym obie metody pozostają dostępne równolegle.

Zakres prac nad modelami obejmuje wydzielenie heurystyki jako nazwanego modelu bazowego za wymienną granicą programistyczną, określenie zdarzenia docelowego wraz ze sposobem jego etykietowania, uruchomienie rejestrowania danych uczących oraz dopasowanie i wdrożenie modelu predykcji opóźnień wraz z jego oceną. Ponieważ dane gromadzone przez sam system narastają dopiero wraz z jego eksploatacją, model dopasowano do publicznego zbioru projektów prowadzonych w metodyce zwinnej.

Poza zakresem pozostaje wykazanie przewagi modelu uczonego nad heurystyką w warunkach konkretnego wdrożenia, ponieważ wymaga to danych o historii dłuższej niż czas realizacji pracy. Z tego powodu obie metody utrzymano w stanie jednocześnie uruchamialnym, a heurystyka pozostała metodą domyślną.

Zakres obejmuje ponadto sformułowanie przydziału zadań jako zadania optymalizacji z ograniczeniami zasobowymi oraz wdrożenie rozwiązującego je solvera obok dotychczasowej heurystyki zachłannej, a także przygotowanie alternatywnej priorytetyzacji backlogu z udziałem modelu językowego.

Priorytetyzacja backlogu uwzględnia wartość biznesową, zależności, złożoność i ocenę ryzyka. Harmonogramowanie obejmuje analizę grafu zależności, wyznaczanie terminów i ścieżki krytycznej oraz wspomaganie przydziału nieprzypisanych zadań. W pracy analizowane są również ograniczenia przyjętego modelu czasu, pojemności zespołu i dopasowania kompetencji.

Estymacja kosztów obejmuje integrację z API Anthropic, przygotowanie kontekstu na podstawie backlogu i stawek zespołu, odbiór ustrukturyzowanej odpowiedzi oraz prezentację wariantów kosztowych, założeń i rekomendacji. Zakres nie obejmuje trenowania dużego modelu językowego od podstaw.

Ocena rozwiązania obejmuje testy oprogramowania oraz eksperymenty dotyczące mechanizmów analitycznych. Wnioski z badań odnoszą się do wykorzystanych danych, scenariuszy i założeń. Ocena wpływu systemu na długoterminowe wyniki rzeczywistych organizacji pozostaje poza zakresem pracy.

\section{Teza pracy}

Na potrzeby pracy przyjęto następującą tezę:

\textbf{Integracja danych projektowych z metodami analizy ryzyka, wielokryterialnej priorytetyzacji, algorytmicznego przydziału zadań i estymacji kosztów z wykorzystaniem dużego modelu językowego umożliwia poprawę jakości wsparcia decyzji planistycznych względem przyjętych metod bazowych w określonych scenariuszach projektowych.}

Jakość wsparcia decyzji jest rozpatrywana przez pryzmat mierzalnych właściwości wyników. Obejmuje trafność identyfikowania opóźnień, zgodność kolejności zadań z kryteriami priorytetyzacji, rozkład obciążenia zespołu i spełnienie ograniczeń przydziału, a także błąd estymacji kosztów względem przyjętych wartości referencyjnych. Dodatkowo analizowana jest możliwość prześledzenia założeń i czynników wpływających na rekomendację.

Weryfikacja tezy wymaga porównania metod przy zachowaniu zgodnych danych wejściowych i warunków eksperymentu. Samo wygenerowanie prognozy lub rekomendacji nie stanowi potwierdzenia jej jakości. Konieczne jest zbadanie, czy uzyskany wynik przewyższa odpowiedni punkt odniesienia oraz w jakich przypadkach zastosowane podejście nie przynosi poprawy.

\section{Zastosowane metody badawcze}

Metodyka pracy obejmuje analizę literatury i istniejących rozwiązań, projektowanie artefaktu informatycznego, testowanie oprogramowania oraz eksperymenty porównawcze. Część eksperymentalna została zrealizowana w odniesieniu do predykcji opóźnień; dla pozostałych mechanizmów przedstawiono zaplanowaną metodykę, a miejsca te oznaczono jawnie, aby opis procedury nie był mylony z potwierdzeniem jej wykonania.

Analiza literatury służy uporządkowaniu zagadnień zarządzania projektami, oceny ryzyka, harmonogramowania, uczenia maszynowego i zastosowania modeli językowych, a analiza porównawcza narzędzi — określeniu zakresu funkcji istniejących systemów i uzasadnieniu decyzji projektowych. Część inżynierska wykorzystuje iteracyjne projektowanie i implementację, obejmujące modelowanie danych i procesów, określanie granic modułów oraz rozwijanie funkcji użytkowych i analitycznych, wraz z weryfikacją opartą na testach architektury, testach jednostkowych i scenariuszach kontroli dostępu.

Badanie predykcji opóźnień opiera się na danych zawierających cechy dostępne w chwili sporządzania prognozy i późniejszy wynik realizacji zadania. Ich gromadzenie uruchomiono w systemie, ponieważ zbioru takiego nie da się odtworzyć wstecz z danych opisujących stan bieżący, natomiast samo dopasowanie modelu przeprowadzono na zbiorze zewnętrznym obejmującym projekty wielu organizacji. Podział danych ograniczał przenikanie informacji między częścią uczącą a testową i wykorzystywał dwa protokoły: grupowanie według projektu, odpowiadające ocenie przedsięwzięć nieznanych modelowi, oraz podział według czasu. Osobno zbadano zawartość danych pod kątem wielkości opisujących przebieg realizacji, które nie mogą pełnić funkcji cech, oraz oszacowań nadanych po rozpoczęciu sprintu.

Dla pozostałych mechanizmów przedstawiono metodykę bez wyników. Obejmuje ona analizę wrażliwości rankingu priorytetów na zmianę wag i porównanie z prostszymi regułami porządkowania, ocenę przydziału zadań względem zadeklarowanej funkcji celu przy tym samym zakresie pracy oraz badanie estymacji kosztów pod kątem poprawności rachunkowej, wykorzystania przekazanych stawek i powtarzalności wyników. Końcowym etapem jest zestawienie wyników z celami pracy i przyjętą tezą, obejmujące zarówno uzyskane korzyści, jak i przypadki nieskuteczności metod.

\section{Struktura pracy}

Praca obejmuje osiem rozdziałów. Pierwszy przedstawia uzasadnienie wyboru tematu, cel i zakres opracowania, tezę oraz metodykę jej weryfikacji. Drugi omawia podstawy zarządzania projektami i wspomagania decyzji: organizację pracy, estymację, zależności między zadaniami, ocenę ryzyka, priorytetyzację, harmonogramowanie z ograniczeniami zasobowymi oraz zastosowanie uczenia maszynowego i modeli językowych. Trzeci zawiera analizę wybranych narzędzi, wymagania funkcjonalne i niefunkcjonalne, przypadki użycia oraz ograniczenia projektu.

Rozdziały czwarty i piąty dotyczą budowy systemu: architektury, obejmującej podział na moduły i warstwy, organizację danych, granice modeli i kontrolę dostępu, oraz implementacji funkcji zarządzania projektami wraz z integracją aplikacji klienckiej z backendem.

Rozdział szósty koncentruje się na mechanizmach wspomagania decyzji. Opisuje ocenę ryzyka, priorytetyzację backlogu, harmonogramowanie, przydział zadań i estymację kosztów wraz z założeniami i ograniczeniami, a osobne sekcje poświęca pracom nad modelami: wymiennym granicom wraz z mechanizmem gromadzenia danych uczących, doborowi zbioru danych oraz dopasowaniu, ocenie i wdrożeniu modelu predykcji opóźnień, sformułowaniu przydziału zadań jako problemu optymalizacyjnego, a także współpracy z zewnętrznymi modelami językowymi.

Rozdział siódmy obejmuje testowanie i ocenę rozwiązania: wyniki wykonania zestawu testów, metodykę eksperymentów, odniesienie uzyskanych wyników modelu do przyjętych kryteriów, pomocnicze studium na zbiorze odrzuconym oraz dyskusję ograniczeń badań. Rozdział ósmy zawiera podsumowanie realizacji celów, omówienie wkładu własnego i propozycje dalszego rozwoju.

 